\documentclass[11pt]{article}
\usepackage[margin=1in]{geometry}
\usepackage{amsmath,amssymb,amsthm}
\newtheorem{theorem}{Theorem}
\title{A strict-sign obstruction to one prescribed denominator contraction}
\author{}
\date{November 24, 2026}
\begin{document}
\maketitle

\paragraph{Status and scope.}
This is a complete written finite-cube certificate, checked by exact
symbolic and integer arithmetic, but not formalized in Lean. It shows
that a particular simultaneous straight contraction of three
denominators can lose a realizable strict sign pattern, even when all
three affine numerators may change arbitrarily. It does not exclude
another denominator path, another common limiting denominator, or
another three-head representation of the same Boolean function. It is
not a lower bound on unrestricted head complexity.

On four Boolean coordinates put
\[
 B=1+x_1+x_2+x_3+x_4,\qquad D=x_1+2x_2,\qquad E=x_3+2x_4,
 \qquad C_1=B,\quad C_2(t)=B+tD,\quad C_3(t)=B+tE.
\]
For every $t>0$, these denominators are positive on the whole cube and
all slopes are strictly positive. Order the sixteen inputs
lexicographically, with $x_4$ changing fastest, and define
\[
 \sigma=(+,-,+,-,-,+,-,-,+,-,-,+,-,-,+,-).
 \tag{1}
\]

\begin{theorem}
The sign table (1) is realized by three admissible heads with
denominators $C_1,C_2(1),C_3(1)$. For every $0<t\leq1/2$, it is not
realized by any three affine numerators over the prescribed
denominators $C_1,C_2(t),C_3(t)$, even with an arbitrary extra score bias.
\end{theorem}

\begin{proof}
For existence at $t=1$, take the following coefficient vectors in the
order $(1,x_1,x_2,x_3,x_4)$:
\[
 \begin{aligned}
 A_1&=(-151022,-108764,268636,21418,8272),\\
 A_2&=(114799,119774,-297673,-21143,0),\\
 A_3&=(54912,49152,-81584,0,0).
 \end{aligned}
\]
The cleared score $P=A_1C_2(1)C_3(1)+A_2BC_3(1)+A_3BC_2(1)$ has
the exact values
\[
 \begin{split}
 (&18689,-3960,3960,-3960,-3960,3960,-67644,-165836,\\
  &3960,-262917,-3960,3960,-3960,-165836,3960,-3960).
 \end{split}                                                     \tag{2}
\]
They are nonzero and have signs (1). Clearing the three positive
denominators preserves those signs.

For nonexistence define eight rational functions
\[
 \begin{aligned}
 a&=(3t+5)(3t^2-16t-16)/6,\\
 b&=(3t+2)(3t+5)(5t+4)/3,\\
 c&=-2t(3t+5)(9t+8)/9,\\
 d&=-2(3t+5)(9t^3+15t^2-12t-16)/(3(2t+3)),\\
 e&=(3t+5)(6t^2+t-4)/3,\\
 f&=-2(3t+2)(3t+5)(3t^2+12t+8)/(3(t+3)),\\
 g&=2(t+1)(3t+2)(3t+5)/3,
 \qquad h=-2t(t+1)(3t+4).
 \end{aligned}
\]
These names shorten the explicit weight vector
\[
 \lambda_t=(0,a,b,c,a,d,0,e,b,0,f,g,c,e,g,h).             \tag{3}
\]
For every $v\in\{1,x_1,x_2,x_3,x_4\}$, direct expansion gives
\[
 \langle\lambda_t,vC_2(t)C_3(t)\rangle
 =\langle\lambda_t,vBC_3(t)\rangle
 =\langle\lambda_t,vBC_2(t)\rangle=0.                  \tag{4}
\]
All fifteen identities hold symbolically as rational identities in
$t$; the accompanying checker verifies each numerator is exactly zero.

When $0<t\leq1/2$, the denominators displayed above are positive and
\[
 a<0,\quad b>0,\quad c<0,\quad d>0,\quad e<0,
 \quad f<0,\quad g>0,\quad h<0.
\]
For the only three nontrivial sign checks, use
\[
 \begin{aligned}
 3t^2-16t-16&\leq3/4-16<0,\\
 9t^3+15t^2-12t-16&\leq9/8+15/4-16<0,\\
 6t^2+t-4&\leq6/4+1/2-4=-2<0.
 \end{aligned}
\]
Thus every nonzero coordinate of $\lambda_t$ has the corresponding
sign (1), and $\lambda_t$ is not zero. If affine numerators produced
strict signs (1), their cleared score $P_t$ would satisfy
$\langle\lambda_t,P_t\rangle>0$. But (4) and affineness give
$\langle\lambda_t,P_t\rangle=0$, a contradiction.

An extra bias adds a multiple of $BC_2(t)C_3(t)$ to the cleared
score. This term is already in the span annihilated by (4), by taking
the affine multiplier $v=B$ of $C_2(t)C_3(t)$, or equivalently by
absorbing the bias into the first numerator. Hence the argument
includes arbitrary bias.
\end{proof}

\paragraph{Consequence for the proposed normalization.}
The two-head identity in \texttt{two-head-sign-coalescence.tex} cannot
be extended by simply moving every other denominator toward one
reference denominator while freely readjusting affine numerators.
The example excludes an entire interval of positive contraction
parameters and not just preservation of one selected cleared score.
It leaves open whether all strict sign patterns with a given head
count have some different realization whose denominators coalesce.
\end{document}
